% Introduction is unnumbered, following the two specified journal articles.
网络化状态估计通过通信网络将传感器测量传送至远端，为系统监测提供状态信息。
受通信资源限制，测量数据通常不能在每个采样时刻全部发送，间歇观测与自适应采样因而成为状态估计中的重要问题\cite{sinopoli2004,rabi2012}。
随机事件触发机制根据测量与预测之间的偏差决定是否通信，并利用静默事件提供的信息更新估计。
Han等\cite{han2015}提出的随机触发规则保持了条件高斯结构，为分析通信过程与估计精度之间的关系提供了基础。
当传输数据受到虚假数据注入攻击时，接收端估计和后续触发过程可能同时改变。
如何利用通信记录评价统计隐蔽攻击造成的真实估计损伤，是此类系统安全分析需要解决的问题。

针对网络攻击下的检测与安全估计，已有研究给出了攻击辨识和状态重构的基本条件\cite{pasqualetti2013,fawzi2014}。
在损伤分析方面，Mo和Sinopoli\cite{mo2016}、Bai等\cite{bai2017,baisensors2017}以及王彩云和李芳菲\cite{targetwang2021}
分别研究了不同攻击通道与检测约束下的估计性能退化。
Hirzallah和Voulgaris\cite{hirzallah2018}在同时限制攻击幅值和监测输出的条件下，利用线性规划求解线性时不变系统的最坏攻击；
Milo\v{s}evi\'{c}等\cite{milosevic2020}则在有限残差序列的Kullback--Leibler（KL）散度约束下分析期望峰值损伤。
这些工作表明，攻击可造成的损伤取决于攻击权限、资源约束及监测信息。

事件触发机制使上述关系进一步受到发送决策的影响。
Shang等\cite{shang2022}在源端本地滤波架构下，结合接收次数限制和给定事件历史后的创新分布匹配，研究可改变接收事件及替换已发送创新的最坏攻击。
Cheng等\cite{cheng2020}考虑可操纵传输决策和反馈的攻击。
Song和Ye\cite{Paper_Song2022_SEBScheduling}通过篡改预测反馈改变发送概率，在通信率与卡方检测约束下优化平均误差协方差迹。
Guo等\cite{guo2025}将随机事件用于攻击调度，在固定调度、给定信息集和单时刻创新条件矩约束下建立全部可行攻击中的最优性。
此外，Zhou等\cite{zhou2023}研究了区间异常检测器下的历史相关攻击，Xu等\cite{xu2026}分析了不同反馈中心与间歇完整性攻击下的长期稳定性边界。

为提高攻击的可检测性，物理水印及其检测延迟、不确定性问题也受到关注\cite{mo2015,naha2023,targetchen2025}。
Shang等\cite{Paper_Shang2026_SingleDimensional}通过加密源端创新的一个实值线性组合，在事件不变、密文不可篡改及相应谱与投影条件下，
排除了其线性变换与替换噪声模型中的非平凡严格隐蔽攻击。
在载荷完整性尚未获得端到端保证的情况下，还需评价既有可信信息能够约束的剩余损伤。
本文考虑发送事件及实际预测反馈可信、仅已发送载荷可被有界修改的协议，
并用总变差（total variation，TV）描述指定通信记录联合分布与正常分布的差异。
这一处理与完整观测前缀的分布距离分析相联系\cite{khazraei2022perception}，其适用前提是系统能够维持上述可信边界。

对于这一协议，仍有两个问题需要明确。
一是记录的统计差异与真实均方误差之间的定量关系。
载荷修改既引入条件估计偏差，又通过实际预测反馈改变事件历史的概率，后者进一步影响平均名义协方差；
仅计算偏差平方无法完整反映真实损伤。
二是损伤保证所需的记录范围。
远端评估者没有真实已发送测量的副本，其可利用的信息来自事件、反馈及接收载荷，
因而需要分析缩短观察时域或删去历史条件是否仍能保留损伤保证。
已有攻击也需在本文权限下重新检验：Song和Ye的策略读取发送前测量并替换反馈，Guo等的递推使用每拍源端创新历史；
直接移植反向创新规则\cite{shang2022}则不满足有限逐路径峰值，对修改量作径向截断又会改变二阶矩。

针对上述问题，本文在全状态含噪观测、状态转移矩阵可逆的条件下，建立有限时刻的双预算损伤界。
首先分解真实均方误差增量，利用协方差加权的下一静默统计量约束条件偏差，并计入协方差期望的变化；
随后构造满足完整联合日志约束的两拍随机符号攻击，证明固定模型、有限维数、目标时刻及预设有限峰值范围内，
三类记录对应的最坏损伤均具有 $\Theta(\min\{b^2,\varepsilon\})$ 的阶数。
进一步给出一般维数的当拍截止日志反例和标量末事件位反例，分析保证对观察时域与历史信息的依赖。
最后，通过二维耦合系统的数值仿真验证损伤分解、攻击构造的预算可行性以及记录观察时域的作用。
